\clearpage
\onecolumn
\thispagestyle{empty}
\vspace*{\fill}
\begin{center}
	{\color{radblue}\rule{\linewidth}{2pt}}\par
	\vspace{6mm}
	{\Huge\bfseries\color{radblue} Tórax}\par
	\vspace{6mm}
	{\color{radblue}\rule{\linewidth}{2pt}}
\end{center}
\vspace*{\fill}
\clearpage
\twocolumn
\pagestyle{contenido}


%%%%%%%%%%%%%%%%%%%%%%%%%%%%%%%%%%%%%%%%%%%%%%%%%%%%%%%%%%%
%  SECCIÓN: Tórax
%  Archivo: secciones/torax.tex
%  Sin preámbulo ni \begin{document} — lo maneja main.tex
	%%%%%%%%%%%%%%%%%%%%%%%%%%%%%%%%%%%%%%%%%%%%%%%%%%%%%%%%%%%
	
	\section{Tórax}
	
	\subsection{Tórax Frente PA}
	
	\subsubsection*{1. Identificación}
	
	\textbf{Tipo:} Proyección básica estándar. Forma par con la Proyección Lateral de Tórax.
	
	\textbf{Indicaciones clínicas principales:}
	\begin{itemize}
		\item Neumonía, bronquitis e infecciones del tracto respiratorio inferior
		\item Neumotórax (se complementa con proyección en espiración forzada)
		\item Evaluación de silueta cardiaca: cardiomegalia, derrame pericárdico
		\item Derrame pleural y engrosamiento pleural
		\item Lesiones mediastínicas y masas pulmonares
		\item Control post-quirúrgico, post-intubación y seguimiento de patologías crónicas
	\end{itemize}
	
	\subsubsection*{2. Técnica Base}
	
	\textbf{Receptor de imagen:} 35 × 43 cm. Orientación \textbf{transversal} en paciente hiperesténico y esténico; \textbf{longitudinal} en asténico e hipoesténico. Borde superior 4 a 5 cm por encima de los hombros.
	
	\textbf{DFRI:} 180 a 200 cm. Minimiza la magnificación de estructuras mediastínicas por divergencia del haz y optimiza la radioprotección.
	
	\textbf{Técnica:} Ultra homogeneizada (5 o 6 homogenizaciones). Produce escala de grises larga para visualizar simultáneamente trama pulmonar, mediastino y estructuras óseas, minimizando contraste innecesario y dosis al paciente.
	
	\textbf{Factores técnicos:} kVp \textbf{110} --- mAs \textbf{2,5}. Espesor medido a los \textbf{2/3 del tórax} (plano de los pechos). Fórmula base: mAs = 10 × 3,24 × 3,5 $\approx$ 114, dividido por 2 en cada homogenización. La pérdida de mAs se compensa con Bucky y mayor DFRI.
	
	\textbf{Bucky:} Bucky mural (vertical) con grilla antidifusora.
	
	\subsubsection*{3. Posicionamiento}
	
	\paragraph{Paciente}
	Bipedestación o sedestación. Plano medio sagital perpendicular al receptor y centrado en su línea media. Siempre en PA: el corazón y el mediastino quedan cerca del receptor reduciendo su magnificación; protege tiroides, tejido mamario y esternón (órgano hematopoyético primario). Pies bien apoyados y separados. El ortostatismo desciende el diafragma por gravedad, evita la ingurgitación vascular pulmonar y permite demostrar niveles hidroaéreos.
	
	\paragraph{Región de interés}
	Hombros descendidos y en contacto firme con el receptor, ambos en el mismo plano transversal, descendiendo las clavículas y despejando los vértices. Brazos en jarra: dorso de las manos sobre las caderas, codos orientados hacia el receptor; esto rota internamente los húmeros y desplaza las escápulas fuera de los campos pulmonares. Mentón extendido levemente hacia arriba para no bloquear vías aéreas ni vértices. En pacientes con mamas voluminosas, pedirle que las eleve y separe lateralmente para despejar las bases pulmonares.
	
	% TODO: \usepackage{graphicx} required
	\begin{figure}[h]
		\centering
		\includegraphics[width=0.7\linewidth]{imagenes/torax position}
		\caption{}
		\label{fig:torax-position}
	\end{figure}
	
	
	\paragraph{Rayo Central}
	Perpendicular al plano del receptor (0°). Punto de entrada: \textbf{T7}, usando el ángulo inferior de las escápulas como reparo anatómico.
	
	\paragraph{Receptor de imagen}
	Centrado en T7. Borde superior 4 a 5 cm por encima de los hombros para incluir los vértices. El plano medio sagital del paciente coincide con la línea media del receptor.
	
	\paragraph{Respiración}
	Inspiración forzada máxima, sostenida durante la exposición. Garantiza la máxima expansión del tórax en los tres ejes y el descenso del diafragma. En sospecha de neumotórax: dos exposiciones, una al final de inspiración forzada y otra al final de espiración forzada. La comparación evidencia aire libre pleural oculto en inspiración; también demuestra movilidad diafragmática, cuerpo extraño y atelectasia.
	
	\subsubsection*{4. Anatomía Visible}
	
	Tráquea aireada en línea media. Ambos campos pulmonares completos desde vértices hasta ángulos costofrénico y cardiofrénico. Cúpulas diafragmáticas. Silueta cardiaca y botón aórtico. Trama vascular pulmonar y árbol bronquial. Costillas, clavículas, escápulas proyectadas fuera de los campos y columna torácica. La trama pulmonar debe ser visible a través de la silueta cardiaca. Los ángulos costofrénico y cardiofrénico deben visualizarse íntegros.
	
	\subsubsection*{5. Criterios de Evaluación}
	
	\textbf{Posicionamiento correcto:} campos pulmonares íntegros incluyendo vértices y ángulos costofrénico y cardiofrénico; \textbf{9 a 10 arcos costales posteriores} (horizontales) y \textbf{6 a 7 arcos costales anteriores} (curvos) visibles por encima de las cúpulas; contornos definidos del corazón y el diafragma; marcas pulmonares desde el hilio a la periferia; marca de lado derecho en la imagen.
	
	\textbf{Ausencia de rotación:} extremidades mediales de ambas clavículas equidistantes de la columna (si la izquierda parece más corta, el paciente está rotado hacia la izquierda); apófisis espinosas centradas en la línea media; tráquea en la línea media.
	
	\textbf{Penetración y contraste:} trama pulmonar visible a través de la silueta cardiaca, confirmando la técnica ultra homogeneizada. Pulmones, mediastino y estructuras óseas diferenciados en la misma imagen.
	
	\subsubsection*{6. Puntos Críticos}
	
	\textbf{Rotación del paciente:} asimetría de clavículas y distorsión de la sombra cardiaca. Corrección: pedir que adelante el hombro más alejado del Bucky y lo presione hasta que ambos queden en contacto firme y parejo. Si una clavícula aparece más alta, indicarle que al inspirar no eleve los hombros y los relaje en un plano horizontal paralelo al piso.
	
	\textbf{Escápulas superpuestas a los campos:} ocurre cuando los brazos no están en jarra o los codos no se orientan hacia adelante. La rotación interna del húmero es la que desplaza las escápulas.
	
	\textbf{Inspiración insuficiente:} menos de 9 costillas posteriores visibles, campos condensados y corazón aparentemente agrandado. Indicador de inspiración correcta: 9--10 costillas posteriores y 6--7 anteriores, con trama pulmonar visible a través del corazón.
	
	\textbf{Corte de ángulos:} pérdida de información diagnóstica crítica; repetir siempre que sea posible.
	
	\textbf{Adaptaciones:} paciente que no puede estar de pie: AP en sedestación o decúbito, la mayor magnificación cardiaca es esperada y debe consignarse. Mamas voluminosas: elevarlas y separarlas. Clavícula congénitamente más alta: no confundir con rotación.
	
	\textbf{Trucos:} espesor medido a los 2/3 del tórax, no en el punto más grueso. La DFRI larga compensa la pérdida de mAs de las homogenizaciones. Si no se ve trama a través del corazón, el kVp o las homogenizaciones fueron insuficientes. La PA magnifica menos el corazón que la AP, relevante para medir el índice cardiotorácico.
	
	
	\clearpage
	
	\begin{figure*}[h]
		\centering
		\includegraphics[width=0.7\linewidth]{imagenes/image}
		\caption[]{Torax PA}
		\label{fig:frentedetoraxPA}
	\end{figure*}
	
\clearpage
	%----------------------------------------------------------
	%  Para agregar otro enfoque duplicá el bloque
	%  desde \subsection hasta \textbf{Trucos}
	%----------------------------------------------------------


%%%%%%%%%%%%%%%%%%%%%%%%%%%%%%%%%%%%%%%%%%%%%%%%%%%%%%%%%%%
%  SECCIÓN: Tórax
%  Archivo: secciones/torax.tex
%  Sin preámbulo ni \begin{document} — lo maneja main.tex
	%%%%%%%%%%%%%%%%%%%%%%%%%%%%%%%%%%%%%%%%%%%%%%%%%%%%%%%%%%%
	
	\section{Tórax}
	
	\subsection{Proyección Lateral de Tórax (D o I)}
	
	\subsubsection*{1. Identificación}
	
	\textbf{Tipo:} Proyección básica. Forma par con la proyección posteroanterior (PA) de tórax.
	
	\textbf{Indicaciones clínicas principales:}
	\begin{itemize}
		\item Demostración de patologías situadas por detrás del corazón, grandes vasos y el esternón.
		\item Demostración de fisuras interlobulares y diferenciación de lóbulos pulmonares.
		\item Localización de lesiones pulmonares (complementa la proyección frontal).
		\item Hemotórax (presencia de sangre en cavidad pleural).
		\item Neumotórax (pérdida de trama pulmonar).
	\end{itemize}
	
	\subsubsection*{2. Técnica Base}
	
	\textbf{Receptor de imagen:} 35 × 43 cm. Orientación \textbf{longitudinal (vertical)}, con el borde superior ubicado 4 a 5 cm por encima de los hombros.
	
	\textbf{DFRI:} 180 a 200 cm. Distancia aumentada para minimizar la magnificación de las estructuras mediastínicas (reducir la divergencia del haz) y mejorar el detalle óseo de las estructuras torácicas. Contribuye además a la radioprotección del paciente.
	
	\textbf{Técnica:} Ultra homogeneizada (5 o 6 homogenizaciones). Permite obtener una escala de grises larga, reduciendo el contraste y posibilitando la visualización simultánea del parénquima pulmonar, el mediastino y las estructuras óseas, incluyendo la trama pulmonar detrás de las costillas y el corazón.
	
	\textbf{Factores técnicos:} kVp \textbf{igual al utilizado en la proyección frontal} --- mAs \textbf{= mAs del frente × 4} (factor 3,5 estándar; factor 4 en pacientes obesos). El factor de multiplicación compensa el mayor espesor que debe atravesar el rayo central en la proyección lateral.
	
	\textbf{Bucky:} Bucky mural con grilla. Obligatorio dado el espesor de la región y la técnica de alta tensión utilizada.
	
	\subsubsection*{3. Posicionamiento}
	
	\paragraph{Paciente}
	Paciente en bipedestación o sedestación. Plano medio sagital (PMS) paralelo al receptor de imagen. El peso corporal se distribuye de manera equitativa entre ambas piernas. La posición vertical favorece el descenso del diafragma por acción de la gravedad, permitiendo mejor visualización de los campos pulmonares y los ángulos costofrénicos, facilita la demostración de niveles hidroaéreos y evita la ingurgitación de los vasos sanguíneos pulmonares.
	
	\paragraph{Región de interés}
	Se coloca el hombro adyacente al receptor apoyado sobre el mismo. Se indica al paciente que levante los brazos por encima de la cabeza, agarrándose los codos entre sí o sosteniéndose de un soporte, para evitar que las partes blandas de los miembros superiores se superpongan a los campos pulmonares. En pacientes inestables se coloca un soporte auxiliar. Se indica elevar el mentón y mirar hacia adelante. El plano medio coronal se centra en la línea media del receptor para evitar rotaciones.
	
	\paragraph{Rayo Central}
	Perpendicular al plano del receptor. Punto de entrada: \textbf{plano medio coronal a la altura de T7 (equivalente al ángulo inferior de las escápulas)}.
	
	\paragraph{Receptor de imagen}
	Centrado a nivel de T7. Borde superior 4 a 5 cm por encima de los hombros, para incluir los vértices pulmonares. La región abarcada comprende desde las clavículas hasta los ángulos costofrénicos, incluyendo ambos campos pulmonares y el mediastino completo.
	
	\paragraph{Respiración}
	Exposición al final de una inspiración forzada (máxima). Garantiza la máxima expansión del tórax en los tres ejes, permitiendo correcta visualización de los campos pulmonares y los ángulos costofrénicos. Luego del disparo se indica al paciente que retome la respiración con normalidad.
	
	\begin{figure}[h]
		\centering
		\includegraphics[width=0.7\linewidth]{imagenes/posicion torax perfil}
		\caption{}
		\label{fig:posicion-torax-perfil}
	\end{figure}
	
	
	
	
	
	\subsubsection*{4. Anatomía Visible}
	
	Corazón y aorta. Fisuras interlobulares (diferenciación de lóbulos pulmonares). Campos pulmonares en toda su extensión vertical. Mediastino anterior y posterior. Esternón de perfil. Cuerpos vertebrales torácicos con espacios intervertebrales y agujeros intervertebrales. Hilios pulmonares (ubicados aproximadamente en el centro de la radiografía). Ángulos costofrénicos. Hemidiafragmas. Arcos costales posteriores superpuestos a la columna vertebral.
	
	\subsubsection*{5. Criterios de Evaluación}
	
	\textbf{Posicionamiento correcto:} Esternón en posición de perfil verdadero. Eje largo de los campos pulmonares demostrado en posición vertical, sin inclinación anterior ni posterior. Ángulos costofrénicos y porción inferior de los vértices pulmonares incluidos. Brazos y partes blandas de los miembros superiores sin superponerse a la porción superior de los campos pulmonares. Hilios ubicados en el centro aproximado de la imagen.
	
	\textbf{Ausencia de rotación:} Arcos costales posteriores superpuestos a los cuerpos vertebrales. El lado alejado del receptor puede aparecer mínimamente por detrás del lado apoyado por efecto de la divergencia del haz (hallazgo normal esperado). Bordes posteriores de las vértebras visibles como una línea única. Espacios intervertebrales torácicos abiertos y agujeros intervertebrales visualizables (excepto en pacientes con escoliosis). Contornos nítidos del corazón y el diafragma.
	
	\textbf{Penetración y contraste:} Escala de grises larga que permita visualizar simultáneamente el parénquima pulmonar, el mediastino y las estructuras óseas. Penetración adecuada de los campos pulmonares y del corazón.
	
	\subsubsection*{6. Puntos Críticos}
	
	\textbf{Rotación del paciente:} Distorsiona la sombra cardíaca y compromete la evaluación del mediastino. Corrección: recolocar los hombros de modo que el PMS quede paralelo al receptor y el plano medio coronal perpendicular al mismo.
	
	\textbf{Brazos no elevados correctamente:} Las partes blandas de los miembros superiores se superponen a los vértices pulmonares, enmascarando patología. Corrección: indicar al paciente que lleve los brazos bien por encima de la cabeza, agarrándose los codos o sosteniéndose de un soporte.
	
	\textbf{Adaptaciones:} En pacientes con trauma o imposibilidad de bipedestación se realiza en sedestación o en decúbito lateral con rayo horizontal. En pacientes inestables se utiliza soporte auxiliar (gotero, barra) para mantener la posición de los brazos.
	
	\textbf{Trucos:} La proyección lateral izquierda es el estándar de referencia: el corazón queda más próximo al receptor, minimizando su magnificación y permitiendo mejor evaluación de la silueta cardíaca. Solo se realiza lateral derecha cuando está específicamente indicado (por ejemplo, para evaluar el pulmón derecho en detalle). El factor de multiplicación de mAs (×4) se aplica siempre respecto a los mAs de la proyección frontal, no de forma independiente.
	
	%----------------------------------------------------------
	%  Para agregar otro enfoque duplicá el bloque
	%  desde \subsection hasta \textbf{Trucos}
	%----------------------------------------------------------


\begin{figure*}[h]
	\centering
	\includegraphics[width=0.7\linewidth]{imagenes/torax perfil rx}
	\caption{}
	\label{fig:torax-perfil-rx}
\end{figure*}




\clearpage
	
%%%%%%%%%%%%%%%%%%%%%%%%%%%%%%%%%%%%%%%%%%%%%%%%%%%%%%%%%%%
%  SECCIÓN: Tórax
%  Archivo: secciones/torax.tex
%  Sin preámbulo ni \begin{document} — lo maneja main.tex
	%%%%%%%%%%%%%%%%%%%%%%%%%%%%%%%%%%%%%%%%%%%%%%%%%%%%%%%%%%%
	
	\section{Tórax}
	
	\subsection{Tórax Portátil (Proyección AP)}
	
	\subsubsection*{1. Identificación}
	
	\textbf{Tipo:} Proyección de emergencia/adaptada. Sustituto de la proyección PA estándar cuando el paciente no puede ser trasladado.
	
	\textbf{Indicaciones clínicas principales:}
	\begin{itemize}
		\item Pacientes críticos o no colaborativos en UCI o sala de internación.
		\item Control de posición de dispositivos médicos (sondas, catéteres, tubo endotraqueal).
		\item Seguimiento de patología torácica en pacientes con imposibilidad de bipedestación.
		\item Sospecha de derrame pleural, neumotórax o consolidación en contexto de urgencia.
	\end{itemize}
	
	\subsubsection*{2. Técnica Base}
	
	\textbf{Receptor de imagen:} 35 × 43 cm. Orientación \textbf{transversal} en pacientes hiperesténicos y esténicos; \textbf{longitudinal} en pacientes asténicos e hipoesténicos, para incluir correctamente el eje longitudinal del tórax según el hábito corporal. Borde superior 4 a 5 cm por encima de los hombros.
	
	\textbf{DFRI:} La mayor distancia que permita el equipo y la situación clínica. \textbf{Mínima de 100 cm.} A menor distancia se incrementa la magnificación del mediastino; con la distancia habitual de trabajo portátil se genera una magnificación aproximada del 15\% respecto a la proyección PA estándar.
	
	\textbf{Técnica:} Ajustada a la distancia real de trabajo y a la ausencia de Bucky. Se reduce el kVp respecto a la técnica estándar para disminuir la radiación secundaria por efecto Compton, compensando la falta de grilla y mejorando el contraste de la imagen.
	
	\textbf{Factores técnicos:} kVp \textbf{reducido respecto a la técnica estándar} --- mAs \textbf{ajustado a la distancia utilizada}. Sin método de medición de espesor estandarizado; se ajusta según las condiciones del paciente y el equipo disponible.
	
	\textbf{Bucky:} No disponible. La ausencia de Bucky y de grilla es una limitación inherente a la técnica portátil; genera mayor ruido y menor contraste respecto a la técnica estándar.
	
	\subsubsection*{3. Posicionamiento}
	
	\paragraph{Paciente}
	Paciente lo más erguido posible, idealmente a 90° en la camilla (sedestación o semisentado). La posición vertical favorece la detección de derrames pleurales, neumotórax y niveles hidroaéreos. En posición supina o con escasa elevación estas patologías son más difíciles de detectar por pérdida del efecto gravitacional.

	\paragraph{Región de interés}
	Se indica rodar los hombros del paciente hacia adelante, rotando los brazos medial o internamente, para alejar las escápulas de los campos pulmonares. Se centra al paciente respecto al rayo central y al receptor, verificando la alineación visionando al paciente desde la parte superior, cerca del tubo.


	\paragraph{Rayo Central}
	Perpendicular al receptor y al esternón. Si es necesario, se angula caudalmente (aproximadamente ±5°) hasta lograr la perpendicularidad respecto al eje longitudinal del esternón, para impedir que las clavículas oculten los vértices pulmonares. Punto de entrada: \textbf{T7, aproximadamente 8 a 10 cm por debajo de la escotadura yugular}.
	
	\paragraph{Receptor de imagen}
	Colocado detrás del tórax del paciente, centrado a nivel de T7. Borde superior 4 a 5 cm por encima de los hombros. Orientación adaptada al hábito corporal del paciente (ver sección Técnica Base).
	
	\paragraph{Respiración}
	Idealmente suspendida al final de una inspiración completa. Dado que muchos pacientes no pueden colaborar plenamente, se recomienda visualizar el tórax del paciente y disparar en el momento de mayor volumen respiratorio visible.
	
	
	\begin{figure}[h]
		\centering
		\includegraphics[width=0.7\linewidth]{imagenes/torax portatil}
		\caption{}
		\label{fig:torax-portatil}
	\end{figure}
	
	
	\subsubsection*{4. Anatomía Visible}
	
	Campos pulmonares. Silueta cardíaca (con magnificación inherente a la técnica). Mediastino. Estructuras óseas torácicas (costillas, clavículas, esternón proyectado). Hemidiafragmas y ángulos costofrénicos (con limitaciones en posición supina). Dispositivos médicos en caso de estar presentes.
	
	\subsubsection*{5. Criterios de Evaluación}
	
	\textbf{Posicionamiento correcto:} Escápulas alejadas de los campos pulmonares. Tórax centrado respecto al receptor. Inclusión de ambos ángulos costofrénicos y vértices pulmonares. Borde superior del receptor 4 a 5 cm por encima de los hombros.
	
	\textbf{Ausencia de rotación:} Simetría de las articulaciones esternoclaviculares respecto a la columna vertebral. Columna vertebral centrada en la imagen.
	
	\textbf{Penetración y contraste:} Se deben mantener los mismos criterios de valor diagnóstico que en la radiografía estándar. Se aceptan tres variaciones inherentes a la técnica: (1) silueta cardíaca magnificada por el menor DFRI y el aumento de la distancia objeto-receptor; (2) posible ocultamiento de la trama vascular por derrame pleural; (3) visualización de solo 8 a 9 costillas posteriores por inspiración incompleta, con pulmones de mayor densidad aparente al no estar completamente ventilados.
	
	\subsubsection*{6. Puntos Críticos}
	
	\textbf{Magnificación del mediastino:} Consecuencia directa del menor DFRI y de la proyección AP. No es un error técnico sino una limitación estructural de la técnica; debe tenerse en cuenta al evaluar el índice cardiotorácico.
	
	\textbf{Ausencia de Bucky:} Genera mayor ruido y menor contraste. Se compensa reduciendo el kVp para disminuir la radiación secundaria por efecto Compton.
	
	\textbf{Adaptaciones:} En pacientes en decúbito supino estricto se realiza la proyección en esa posición, asumiendo las limitaciones diagnósticas para derrame pleural, neumotórax y niveles hidroaéreos. En pacientes con múltiples dispositivos médicos se solicita ayuda al personal del área para posicionar el receptor sin traccionar ni desconectar los dispositivos.
	
	\textbf{Trucos:} La calidad de imagen de una radiografía portátil nunca debe aceptarse como inferior por el solo hecho de serlo; el objetivo sigue siendo alcanzar los mismos criterios diagnósticos de la proyección estándar dentro de las limitaciones del contexto. Visualizar al paciente desde arriba (cerca del tubo) antes de disparar permite verificar la alineación sin movilizarlo.
	
	
	\begin{figure*}[h]
		\centering
		\includegraphics[width=0.7\linewidth]{imagenes/torax portatil 2}
		\caption{}
		\label{fig:torax-portatil-2}
	\end{figure*}
	
	\clearpage
	
	%----------------------------------------------------------
	%  Para agregar otro enfoque duplicá el bloque
	%  desde \subsection hasta \textbf{Trucos}
	%----------------------------------------------------------
	
	
	%%%%%%%%%%%%%%%%%%%%%%%%%%%%%%%%%%%%%%%%%%%%%%%%%%%%%%%%%%%
	%  SECCIÓN: Tórax
	%  Archivo: secciones/torax.tex
	%  Sin preámbulo ni \begin{document} — lo maneja main.tex
		%%%%%%%%%%%%%%%%%%%%%%%%%%%%%%%%%%%%%%%%%%%%%%%%%%%%%%%%%%%
		
		\section{Tórax}
		
		\subsection{Vértices Pulmonares (Método de Lindblom / Posición Lordótica)}
		
		\subsubsection*{1. Identificación}
		
		\textbf{Tipo:} Proyección complementaria. Debe realizarse luego de la proyección PA y perfil de tórax convencionales.
		
		\textbf{Indicaciones clínicas principales:}
		\begin{itemize}
			\item Tuberculosis pulmonar (bacilo de Koch) y sospecha de focos apicales
			\item Nódulos, masas o tumores en vértices pulmonares
			\item Casquetes apicales y calcificaciones en región de los vértices
			\item Descartar patología pulmonar superior no visible en PA convencional por superposición clavicular
		\end{itemize}
		
		\subsubsection*{2. Técnica Base}
		
		\textbf{Receptor de imagen:} 35 × 43 cm. Orientación \textbf{longitudinal}. Borde superior del chasis ubicado 7--8 cm por encima de los hombros.
		
		\textbf{DFRI:} 180 cm. Igual que en tórax convencional, para reducir la magnificación del corazón y estructuras mediastínicas.
		
		\textbf{Técnica:} Ultra homogeneizada, igual que en tórax PA. Debe ser \textbf{ligeramente mayor} que en la PA convencional, dado el mayor espesor del tórax en posición oblicuada, que requiere mayor penetración para visualizar tejido pulmonar retroóseo.
		
		\textbf{Factores técnicos:} kVp \textbf{similar a tórax PA, levemente aumentado} --- mAs \textbf{levemente aumentado respecto a PA}. Espesor medido con el tórax en posición lordótica.
		
		\textbf{Bucky:} Bucky de pared con grilla, dado el espesor de la región.
		
		\subsubsection*{3. Posicionamiento}
		
		\paragraph{Paciente}
		Bipedestación a unos 30 cm del receptor. Plano medio sagital perpendicular al receptor. Posición AP: el paciente se inclina lentamente hacia atrás, apoyando hombros, cuello y nuca contra el Bucky de pared, manteniendo los pies fijos en el piso.
		
		\paragraph{Región de interés}
		Ambas manos sobre las caderas con las palmas hacia afuera y los hombros proyectados hacia adelante. El dorso de las escápulas queda en contacto con el receptor. La hiperextensión del tórax proyecta las clavículas hacia arriba, liberando los vértices pulmonares de su superposición.
		
		\paragraph{Rayo Central}
		\textbf{Opción preferida (paciente apto):} Perpendicular al plano del receptor. Punto de entrada: \textbf{zona media del esternón}.\\
		\textbf{Opción alternativa (paciente no logra lordosis):} Angulado \textbf{15--20° en sentido cefálico}, centrado en el \textbf{manubrio esternal}. Regla general: \textbf{siempre evitar angular cuando sea posible}.
		
		\paragraph{Receptor de imagen}
		Centrado en los vértices pulmonares. Borde superior 7--8 cm por encima de los hombros. Colimar \textbf{exclusivamente a los vértices pulmonares}, sin necesidad de incluir la totalidad de los campos pulmonares.
		
		\paragraph{Respiración}
		Inspiración máxima. Expande los pulmones y maximiza la visualización del tejido pulmonar apical.
		
		\begin{figure}[h]
			\centering
			\includegraphics[width=0.7\linewidth]{imagenes/torax vertices}
			\caption{}
			\label{fig:torax-vertices}
		\end{figure}
		
		
		\begin{figure}[h]
			\centering
			\includegraphics[width=0.7\linewidth]{imagenes/torax vertices 2}
			\caption[]{Variante}
			\label{fig:torax-vertices-2}
		\end{figure}
		
		
		\subsubsection*{4. Anatomía Visible}
		
		Vértices pulmonares bilaterales libres de superposición clavicular. Clavículas proyectadas por encima de los vértices en disposición horizontal. Primeras y segundas costillas con sus extremos mediales superpuestos a los extremos mediales de las clavículas. Segmentos anterior y posterior de las costillas superpuestos y distorsionados por la oblicuación del tórax. Región apical del parénquima pulmonar despejada para evaluar nódulos, infiltrados o calcificaciones.
		
		\subsubsection*{5. Criterios de Evaluación}
		
		\textbf{Posicionamiento correcto:} Clavículas situadas por encima de los vértices pulmonares, sin superponerse sobre el tejido pulmonar apical.
		
		\textbf{Ausencia de rotación:} Extremos mediales de ambas clavículas equidistantes respecto a la columna vertebral. Clavículas en disposición horizontal y simétrica.
		
		\textbf{Penetración y contraste:} Técnica ultra homogeneizada que permita visualizar el tejido pulmonar a través de la densidad ósea superpuesta. Los segmentos anteriores y posteriores de las costillas deben verse distorsionados y superpuestos, confirmando la correcta posición lordótica.
		
		\subsubsection*{6. Puntos Críticos}
		
		\textbf{Paciente no logra la posición lordótica:} Las clavículas permanecen superpuestas a los vértices. Corrección: angular el RC 15--20° en sentido cefálico, centrando en el manubrio. Introduce distorsión geométrica de la región.
		
		\textbf{Colimación excesiva:} Incluir la totalidad de los campos pulmonares implica irradiación innecesaria. Al ser proyección complementaria, la colimación debe restringirse a los vértices.
		
		\textbf{Adaptaciones:} En pacientes con dolor torácico agudo o limitación de movilidad, optar directamente por la angulación del RC (15--20° cefálico) sin forzar la posición lordótica.
		
		\textbf{Trucos:} Regla práctica de angulación: en AP (paciente frente al tubo) → ángulo craneal para elevar las clavículas; en PA → ángulo caudal para el mismo efecto. La técnica debe aumentarse respecto a la PA convencional por el mayor espesor efectivo del tórax oblicuado.
		
		
		
		\begin{figure*}[h]
			\centering
			\includegraphics[width=0.7\linewidth]{imagenes/torax vertice 3}
			\caption{}
			\label{fig:torax-vertice-3}
		\end{figure*}
		
		\clearpage
		
		
		%----------------------------------------------------------
		%  Para agregar otro enfoque duplicá el bloque
		%  desde \subsection hasta \textbf{Trucos}
		%----------------------------------------------------------
	
	%%%%%%%%%%%%%%%%%%%%%%%%%%%%%%%%%%%%%%%%%%%%%%%%%%%%%%%%%%%
	%  SECCIÓN: Tórax
	%  Archivo: secciones/torax.tex
	%  Sin preámbulo ni \begin{document} — lo maneja main.tex
		%%%%%%%%%%%%%%%%%%%%%%%%%%%%%%%%%%%%%%%%%%%%%%%%%%%%%%%%%%%
		
		\section{Tórax}
		
		\subsection{Oblicuas Anteriores de Corazón y Grandes Vasos (OAD / OAI)}
		
		\subsubsection*{1. Identificación}
		
		\textbf{Tipo:} Proyecciones complementarias. Forman par entre sí: la OAD y la OAI se
		solicitan conjuntamente según la indicación clínica. \textbf{Solo se realizan oblicuas
			anteriores} para corazón y grandes vasos; las oblicuas posteriores están contraindicadas
		(mayor magnificación y exposición frontal del paciente al haz primario).
		
		\textbf{Nomenclatura:} Se nombran según el lado del cuerpo que se apoya contra el RI.
		
		  \begin{tcolorbox}[colback=red!5!white, colframe=red!80!black, fonttitle=\bfseries, title=Importante]
			\begin{itemize}
				\item \textbf{OAD} (Oblicua Anterior Derecha): hombro anterior \textbf{derecho} apoyado en el RI.
				\item \textbf{OAI} (Oblicua Anterior Izquierda): hombro anterior \textbf{izquierdo} apoyado en el RI.
			\end{itemize}
		\end{tcolorbox}
		
		\textbf{Principio de visualización:} El lado de interés diagnóstico es el más \textbf{alejado}
		del RI. Para ver la hemitórax izquierda (cayado aórtico, campo pulmonar izquierdo) → OAI;
		para ver la hemitórax derecha → OAD.
		
		\textbf{Indicaciones clínicas principales:}
		\begin{itemize}
			\item Bloque quirúrgico: control de stent, malla endovascular o aorta con contraste.
			\item Esofagograma y estudio gastroduodenal con contraste baritado (desplegar el esófago retrocardiaco).
			\item Valoración del cayado aórtico y aorta descendente (preferentemente OAI).
			\item Estudio del mediastino: separar la columna vertebral de las estructuras mediastínicas
			superpuestas en el frente PA estándar.
			\item Evaluación del campo pulmonar completo libre de superposición cardiaca o columnar.
		\end{itemize}
		
		\subsubsection*{2. Técnica Base}
		
		\textbf{Receptor de imagen:} 35 × 43 cm. Orientación \textbf{longitudinal} en todo tipo de paciente,
		para incluir ambos pulmones desde los vértices hasta los ángulos costofrénicos.
		
		\textbf{DFRI:} \textbf{180 cm}. La distancia foco-receptor incrementada reduce la magnificación
		de las estructuras mediastínicas, que en proyección PA quedan alejadas del RI, mejorando la
		nitidez geométrica y la fidelidad de tamaño.
		
		\textbf{Técnica:} Ultra-homogeneizada con \textbf{alto kVp} (escala larga de grises). Permite
		visualizar simultáneamente el parénquima pulmonar y el mediastino, que de otro modo quedarían
		enmascarados por la superposición ósea y la densidad cardiaca. Variantes según paciente:
		\begin{itemize}
			\item \textit{Paciente obeso:} técnica intermedia; mayor radiación dispersa, puede requerirse
			incremento controlado de mAs.
			\item \textit{Control de stent o malla metálica:} kVp muy alto para ``atravesar'' la densidad
			cardiaca y revelar el metal con claridad, priorizando transparencia.
		\end{itemize}
		
		\textbf{Factores técnicos:} kVp \textbf{alto} --- mAs según hábito corporal.
		Espesor medido a nivel del tórax en su diámetro AP.
		
		\textbf{Bucky:} Sí. Bucky mural (vertical) con grilla, dado el espesor del tórax y la
		producción significativa de radiación dispersa.
		
		\subsubsection*{3. Posicionamiento}
		
		\paragraph{Paciente}
		Bipedestación (de pie) o, si no es posible, sedestación erguida. Posición PA: el tórax
		enfrenta el RI y la espalda queda dirigida hacia el tubo, lo que protege las estructuras
		radiosensibles anteriores (radioprotección activa). El mentón elevado, mirando hacia adelante,
		sin rotar la cabeza.
		
		\begin{figure}[h]
			\centering
			\includegraphics[width=0.7\linewidth]{imagenes/grandes vasos OAD}
			\caption[]{OAD grandes vasos}
			\label{fig:grandes-vasos-oad}
		\end{figure}
		
		\begin{figure}[h]
			\centering
			\includegraphics[width=0.7\linewidth]{imagenes/grandes vasos OAI}
			\caption[]{OAI grandes vasos}
			\label{fig:grandes-vasos-oai}
		\end{figure}
		
		
		\paragraph{Región de interés}
		Rotación del tórax sobre su eje longitudinal alejando el lado de interés del RI:
		\textbf{45°} para estudio del \textbf{pulmón} y \textbf{60°} para estudio del
		\textbf{corazón y grandes vasos}, medidos desde el plano frontal.
		\begin{itemize}
			\item \textit{OAD:} se aleja el hombro izquierdo; el hombro anterior derecho queda apoyado en el RI.
			\item \textit{OAI:} se aleja el hombro derecho; el hombro anterior izquierdo queda apoyado en el RI.
		\end{itemize}
		Posición de los miembros superiores: del \textit{lado apoyado}, la mano se coloca sobre la
		cadera con la palma hacia afuera; del \textit{lado alejado}, el paciente eleva el brazo hasta
		la altura del hombro y apoya la mano sobre la cabeza para no tapar el campo pulmonar.
		Variante: ambos codos tomados por encima de la cabeza. Los hombros deben quedar en un mismo
		plano horizontal.
		
		\paragraph{Rayo Central}
		Perpendicular al plano del RI, sin angulación. Punto de entrada:
		\begin{itemize}
			\item \textbf{OAD:} ángulo inferior de la escápula \textbf{izquierda} (lado alejado).
			\item \textbf{OAI:} ángulo inferior de la escápula \textbf{derecha} (lado alejado).
		\end{itemize}
		Referencia alternativa equivalente: \textbf{T7}, aproximadamente 8--10 cm por debajo de la
		escotadura yugular (referenciando C7).
		
		\paragraph{Receptor de imagen}
		Centrado respecto al eje del tórax del paciente. Borde superior situado
		\textbf{4--5 cm por encima de C7} (vértebra prominente). Debe incluir ambos pulmones completos:
		desde los vértices pulmonares hasta los ángulos costofrénicos inferiores, y la totalidad de
		las estructuras mediastínicas (tráquea, cayado aórtico, silueta cardiaca).
		
		\paragraph{Respiración}
		\textbf{Inspiración máxima sostenida} (``Respire todo lo que pueda y aguante el aire'').
		Desciende el diafragma, expande los campos pulmonares y aumenta el contraste aire/tejido blando,
		mejorando la visualización del parénquima y las estructuras vasculares. Control: deben contarse
		\textbf{9--10 arcos costales posteriores} por encima de la cúpula diafragmática (los arcos
		posteriores son más horizontales y más fáciles de identificar). Si solo se observan 8,
		la inspiración fue insuficiente y debe repetirse.
		
		\subsubsection*{4. Anatomía Visible}
		
		\textbf{En OAD:} silueta cardiaca de morfología \textbf{triangular}; espacio aéreo
		retroesternal anterior con forma triangular; cayado aórtico \textit{sin desplegar}
		(superpuesto); tráquea llena de aire y árbol bronquial \textbf{izquierdo} completamente
		desplegado; cámara gástrica \textbf{alejada} de la columna (``posición de esgrimista'');
		campo pulmonar izquierdo en mayor extensión.
		
		\textbf{En OAI:} silueta cardiaca de morfología \textbf{en pera} (ventrículo izquierdo
		proyectado hacia la izquierda inferior); \textbf{cayado aórtico desplegado} con aorta
		ascendente, cayado y aorta descendente separados y nítidos; espacio aéreo anterior de
		morfología \textbf{rectangular}; cámara gástrica \textbf{próxima} a la columna
		(``posición de boxeador''); campo pulmonar derecho en mayor extensión; esófago retrocardiaco
		visible cuando se usa contraste baritado.
		
		En ambas proyecciones: bifurcación traqueal (carina) proyectada \textbf{sobre la silueta
			cardiaca}, nunca sobre la columna vertebral.
		
		\subsubsection*{5. Criterios de Evaluación}
		
		\textbf{Posicionamiento correcto:} OAD → corazón triangular, espacio aéreo anterior
		triangular, cámara gástrica alejada de la columna, árbol bronquial izquierdo desplegado.
		OAI → corazón en pera, cayado aórtico desplegado, espacio aéreo anterior rectangular,
		cámara gástrica próxima a la columna. En ambas: carina sobre la silueta cardiaca y hombros
		en el mismo plano horizontal.
		
		\textbf{Ausencia de rotación:} la columna vertebral no debe superponerse al campo pulmonar
		contralateral (exceso de rotación); el corazón no debe quedar encimado a la columna (falta
		de rotación). La tráquea debe mostrar un perfil aéreo nítido y continuo.
		
		\textbf{Penetración y contraste:} escala larga de grises que permita visualizar
		simultáneamente el parénquima pulmonar y el mediastino. La trama vascular debe ser visible
		hasta la periferia. El corazón no debe ser completamente radiopaco: deben poder evaluarse
		las estructuras retrocardiacas. 9--10 arcos costales posteriores visibles sobre el diafragma
		confirman inspiración adecuada.
		
		\subsubsection*{6. Puntos Críticos}
		
		\textbf{Corazón encimado a la columna (OAI):} falta de rotación; el mediastino no se despliega.
		Corrección: aumentar la rotación hasta 60° y verificar posición de hombros.
		
		\textbf{Columna tapando el campo pulmonar contralateral:} rotación excesiva (demasiado de
		perfil). Corrección: reducir el grado de rotación.
		
		\textbf{Corazón muy radiopaco sin poder atravesarlo:} kVp insuficiente. Corrección: subir kVp
		para obtener escala larga de grises.
		
		\textbf{Imagen con apariencia de parrilla costal:} técnica ósea aplicada erróneamente;
		no se ve la trama pulmonar. Corrección: técnica ultra-homogeneizada con alto kVp.
		
		\textbf{Adaptaciones:} en pacientes con limitación para la bipedestación, realizar en
		sedestación erguida. En bloque quirúrgico (control de stent o malla), priorizar alto kVp
		con escala muy larga para revelar el material metálico a través de la densidad cardiaca.
		
		\textbf{Trucos:}
		\begin{itemize}
			\item Para desplegar el cayado aórtico: siempre OAI. En OAD el cayado queda superpuesto
			y no es valorable.
			\item Nemotecnia: \textit{esgrimista} = OAD (cámara gástrica alejada de la columna);
			\textit{boxeador} = OAI (cámara gástrica próxima a la columna).
			\item Contar arcos \textbf{posteriores} para evaluar la inspiración, no los anteriores;
			son más horizontales y más fáciles de identificar.
			\item Si en un examen preguntan si se realiza oblicua \textit{posterior} para corazón
			→ \textbf{FALSO}: solo se realizan oblicuas anteriores.
		\end{itemize}
		
		\begin{figure*}[t]
			\centering
			\includegraphics[width=0.7\linewidth]{imagenes/grandes vasos RX OAD}
			\caption[]{OAD de grandes vasos}
			\label{fig:grandes-vasos-rx-oad}
		\end{figure*}
		
	\clearpage
	
		\begin{figure*}[t]
			\centering
			\includegraphics[width=0.7\linewidth]{imagenes/grandes vasos RX OAI}
			\caption[]{OAI de grandes vasos}
			\label{fig:grandes-vasos-rx-oai}
		\end{figure*}
		
\clearpage
		
		%----------------------------------------------------------
		%  Para agregar otro enfoque duplicá el bloque
		%  desde \subsection hasta \textbf{Trucos}
		%----------------------------------------------------------
		
		

	
	%%=========================================================
	%%  ENFOQUE 1: Parrilla Costal — Frente (Supra diafragmática)
	%%=========================================================
	\section{Tórax}
	\subsection{Parrilla Costal Frente Supra Diafragmática}
	
	\subsubsection*{1. Identificación}
	
	\textbf{Tipo:} Proyección básica. Primer estudio de la parrilla costal supra diafragmática; forma par con la proyección oblicua costal correspondiente.
	
	\textbf{Indicaciones clínicas principales:}
	\begin{itemize}
		\item Traumatismo de parrilla costal con sospecha de fractura de arcos costales superiores (1.º--10.º)
		\item Dolor costal agudo o crónico de localización superior al diafragma
		\item Evaluación post-protocolo de tórax PA para descartar neumotórax o hemotórax
	\end{itemize}
	
	\textbf{Nota de anamnesis:} Antes de iniciar el estudio verificar: (1) naturaleza y mecanismo del traumatismo; (2) localización exacta del dolor; (3) presencia de dificultad respiratoria; (4) capacidad del paciente para la bipedestación.
	
	\subsubsection*{2. Técnica Base}
	
	\textbf{Receptor de imagen:} 35 $\times$ 43 cm. Orientación \textbf{longitudinal}.
	
	\textbf{DFRI:} 100 cm. Distancia estándar de tórax; reduce la ampliación geométrica y permite incluir la totalidad de la caja torácica superior.
	
	\textbf{Técnica:} \textbf{Homogeneizada o intermedia} (de preferencia según biotipo del paciente). Las costillas supra diafragmáticas están rodeadas de parénquima pulmonar aireado, por lo que no se requiere técnica básica de alta diferenciación de partes blandas.
	
	\textbf{Factores técnicos:} kVp \textbf{65--70}. Un kVp relativamente bajo mantiene el contraste radiológico y permite visualizar las costillas a través del pulmón lleno de aire. Si la lesión se proyecta sobre la silueta cardíaca, incrementar el kVp para obtener mayor penetración en esa región.
	
	\textbf{Bucky:} Bucky mural con grilla.
	
	\subsubsection*{3. Posicionamiento}
	
	\paragraph{Paciente}
	Bipedestación. La posición de pie favorece el descenso gravitacional del diafragma a su posición más baja, ampliando el campo pulmonar y facilitando una inspiración profunda. La orientación del estudio depende de la región a evaluar:
	\begin{itemize}
		\item \textbf{Arcos posteriores}: posición AP (dorso apoyado contra el receptor).
		
		\item \textbf{Arcos anteriores}: posición PA (tórax apoyado contra el receptor; los arcos anteriores quedan más próximos al receptor mejorando el detalle).
	\end{itemize}
	
	\paragraph{Región de interés}
	Plano medio sagital perpendicular al receptor e isocentrado. Se apoya siempre la región que se desea radiografiar para reducir la distancia objeto-receptor y minimizar la ampliación.
	

	
\paragraph{Rayo Central}
Perpendicular al plano del receptor.

\textbf{Bilateral:} centrado en \textbf{T7} (equivalente a la mitad del esternón).

\textbf{Unilateral:} centrado en la \textbf{línea medio clavicular} del hemitórax de interés, a la misma altura.
	
	\paragraph{Receptor de imagen}
	Centrado en T7 para estudios bilaterales, o en la línea medio clavicular ipsilateral para estudios unilaterales. El borde superior del receptor debe superar los ápices pulmonares y el borde inferior incluir hasta por debajo del diafragma.
	
	\paragraph{Respiración}
	\textbf{Inspiración profunda y apnea.} El descenso del diafragma expande los campos pulmonares y proyecta el diafragma por debajo de la 9.ª--10.ª costilla posterior, evitando la superposición sobre los arcos costales superiores. En pacientes con dolor intenso la inspiración puede ser limitada, por lo que en ocasiones solo se visualizan 8--9 costillas posteriores por encima del diafragma.
	
	\subsubsection*{4. Anatomía Visible}
	
	\textbf{Arcos costales posteriores del 1.º al 10.º} por encima del diafragma, proyectados a través del parénquima pulmonar. En la proyección PA se obtiene mejor detalle de los arcos anteriores por su proximidad al receptor. Los arcos costales anteriores presentan una imagen "cortada" en su porción cartilaginosa, ya que el cartílago no calcificado es radiotransparente. La mayor concentración de fracturas se localiza en el \textbf{borde axilar} de los arcos, zona que puede requerir proyecciones oblicuas adicionales para su correcta evaluación.
	
	\subsubsection*{5. Criterios de Evaluación}
	
	\textbf{Posicionamiento correcto:} Visualización de los arcos costales posteriores del 1.º al 10.º por encima del diafragma. En estudios unilaterales, las costillas contralaterales pueden estar incompletas sin pérdida diagnóstica.
	
	\textbf{Ausencia de rotación:} Simetría bilateral de las estructuras torácicas; distancias claviculares equidistantes respecto a la columna vertebral.
	
	\textbf{Penetración y contraste:} Visualización nítida de los arcos costales a través del parénquima pulmonar con contraste adecuado entre hueso y tejido pulmonar aireado.
	
	\subsubsection*{6. Puntos Críticos}
	
	\textbf{Técnica incorrecta (kVp excesivo):} Sobreexposición con pérdida de contraste óseo; los arcos costales se pierden en el fondo del parénquima. Utilizar kVp 65--70 como valor de referencia.
	
	\textbf{Fase respiratoria incorrecta (espiración):} El diafragma asciende y superpone sobre los arcos costales inferiores, reduciendo el número de costillas visibles. Instruir correctamente al paciente para inspiración profunda y apnea.
	
	\textbf{Adaptaciones:} En pacientes con dolor intenso que impide la bipedestación, realizar en decúbito supino (AP acostado); tener en cuenta que el diafragma asciende y puede superponerse sobre los arcos inferiores. Si el dolor impide una inspiración profunda, documentar la limitación y aceptar la imagen con 8--9 costillas visibles.
	
	\textbf{Trucos:} Recordar el \textbf{protocolo obligatorio de dos pasos}: primero radiografía de tórax PA para descartar neumotórax o hemotórax, luego el estudio específico de costillas. La localización del dolor determina si se realiza estudio supra o infra diafragmático.
	
		\begin{figure*}[t]
		\centering
		\includegraphics[width=0.7\linewidth]{imagenes/costilla supra}
		\caption{Arcos costales supra diafragmaticos}
		\label{fig:costilla-supra}
	\end{figure*}
	
\clearpage
	
	%----------------------------------------------------------
	
	%%=========================================================
	%%  ENFOQUE 2: Parrilla Costal — Frente (Infra diafragmática)
	%%=========================================================
	
	\subsection{Parrilla Costal Frente Infra Diafragmática}
	
	\subsubsection*{1. Identificación}
	
	\textbf{Tipo:} Proyección básica. Primer estudio de la parrilla costal infra diafragmática; forma par con la proyección oblicua costal inferior correspondiente.
	
	\textbf{Indicaciones clínicas principales:}
	\begin{itemize}
		\item Traumatismo con sospecha de fractura de arcos costales inferiores (9.ª--12.ª costilla)
		\item Dolor costal agudo o crónico de localización inferior al diafragma
		\item Evaluación post-protocolo de tórax PA para descartar neumotórax o hemotórax
	\end{itemize}
	
	\subsubsection*{2. Técnica Base}
	
	\textbf{Receptor de imagen:} 35 $\times$ 43 cm. Orientación \textbf{longitudinal}.
	
	\textbf{DFRI:} 100 cm. Distancia estándar que permite incluir la región abdominal superior con las costillas inferiores.
	
	\textbf{Técnica:} \textbf{Básica o contrastada} (en pacientes de gran contextura). Las costillas infra diafragmáticas están rodeadas de estructuras abdominales densas y el diafragma muscular, lo que genera mayor dispersión y requiere mayor contraste entre hueso y tejidos blandos.
	
	\textbf{Factores técnicos:} kVp \textbf{70--80}. Un kVp medio garantiza la penetración adecuada a través del diafragma muscular y las estructuras abdominales densas.
	
	\textbf{Bucky:} Bucky mural con grilla.
	
	\subsubsection*{3. Posicionamiento}
	
	\paragraph{Paciente}
	\textbf{Decúbito supino} (preferible al de pie para la región infra diafragmática). En decúbito, el diafragma se eleva a su posición más alta, reduciendo el espesor abdominal efectivo y mejorando la visualización de las costillas inferiores a través de las estructuras abdominales.
	
	\paragraph{Región de interés}
	Plano medio sagital perpendicular al receptor e isocentrado. Se apoya la región posterior para los arcos posteriores.
	
	\paragraph{Rayo Central}
	Perpendicular al plano del receptor. \\
	\textbf{Bilateral:} centrado a nivel de la \textbf{apófisis xifoides}. \\
	\textbf{Unilateral:} centrado en la \textbf{línea medio clavicular} del hemitórax de interés, a la misma altura.
	
	\paragraph{Receptor de imagen}
	Centrado a nivel de la apófisis xifoides para estudios bilaterales. El \textbf{borde inferior de la imagen debe llegar a nivel de las crestas ilíacas} para incluir las costillas 11.ª y 12.ª en su totalidad.
	
	\paragraph{Respiración}
	\textbf{Espiración y apnea.} El ascenso del diafragma lo posiciona a nivel de la 7.ª--8.ª costilla posterior, proporcionando una densidad uniforme sin interfase aire-partes blandas por debajo del diafragma. La ausencia de esta interfase evita artefactos que pueden simular fracturas costales.
	
	\subsubsection*{4. Anatomía Visible}
	
	\textbf{11.ª y 12.ª costillas} (costillas flotantes) por debajo del diafragma, proyectadas a través de las estructuras abdominales. La región abdominal circundante (hígado, bazo, riñones) se visualiza como densidad de partes blandas. El borde inferior de la imagen incluye las crestas ilíacas como referencia de margen inferior.
	
	\subsubsection*{5. Criterios de Evaluación}
	
	\textbf{Posicionamiento correcto:} Visualización de la 11.ª y 12.ª costilla por debajo del diafragma; borde inferior de la imagen a nivel de las crestas ilíacas.
	
	\textbf{Ausencia de rotación:} Simetría bilateral de las estructuras abdominales; columna vertebral centrada en la imagen.
	
	\textbf{Penetración y contraste:} Visualización de los arcos costales inferiores con contraste adecuado a través de la densidad abdominal; sin sobreexposición de partes blandas.
	
	\subsubsection*{6. Puntos Críticos}
	
	\textbf{Fase respiratoria incorrecta (inspiración):} El diafragma desciende y genera una interfase aire-partes blandas que puede simular fracturas costales y dificulta la evaluación. Instruir correctamente al paciente para espiración y apnea.
	
	\textbf{kVp insuficiente:} Subexposición con arcos costales inferiores no visibles a través de la densidad abdominal. Asegurar kVp 70--80 para adecuada penetración.
	
	\textbf{Adaptaciones:} Si el paciente no puede mantenerse en decúbito supino, considerar la bipedestación; tener presente que en esa posición el diafragma desciende y el abdomen se engrosa, pudiendo requerir ajuste de kVp. En pacientes de gran contextura usar técnica contrastada.
	
	\textbf{Trucos:} El cambio de posición (decúbito para infra, bipedestación para supra) responde a la lógica del diafragma: se necesita que esté \textbf{alto} para la región infra y \textbf{bajo} para la región supra. La fase respiratoria es opuesta a la de la proyección supra diafragmática por la misma razón.
	
	
	\begin{figure*}[t]
		\centering
		\includegraphics[width=0.7\linewidth]{imagenes/costillas infra}
		\caption[]{Arcos costales infra diafragmaticos}
		\label{fig:costillas-infra}
	\end{figure*}
	
\clearpage
	%----------------------------------------------------------
	%  Para agregar otro enfoque duplicá el bloque
	%  desde \subsection hasta \textbf{Trucos}
	%----------------------------------------------------------


%%%%%%%%%%%%%%%%%%%%%%%%%%%%%%%%%%%%%%%%%%%%%%%%%%%%%%%%%%%
%  SECCIÓN: Parrilla Costal
%  Archivo: secciones/parrilla_costal.tex
%  Sin preámbulo ni \begin{document} — lo maneja main.tex
	%%%%%%%%%%%%%%%%%%%%%%%%%%%%%%%%%%%%%%%%%%%%%%%%%%%%%%%%%%%
	
	\section{Tórax}
	
	\subsection{Parrilla Costal - Oblicua Anterior PA - Porción Axilar (OAD / OAI)}
	
	\subsubsection*{1. Identificación}
	
	\textbf{Tipo:} Proyección complementaria. Forma par según el lado afectado:
	\textbf{OAI} (hombro anterior izquierdo apoyado) para desplegar la curva axilar
	\textbf{derecha}; \textbf{OAD} (hombro anterior derecho apoyado) para visualizar las
	porciones medias y anteriores del lado \textbf{derecho}. Se realizan siempre luego
	de la radiografía de tórax PA estándar a 1{,}80~m, cuyo objetivo es descartar
	neumotórax, hemotórax o hemoneumotórax y priorizar el hemitórax a estudiar.
	
	\textbf{Indicaciones clínicas principales:}
	\begin{itemize}
		\item Golpe en el borde axilar del tórax.
		\item Lesiones, fisuras o traumatismo costal con dolor localizado post-trauma.
		\item Puntada o dolor en el costado de origen traumático.
	\end{itemize}
	
	\subsubsection*{2. Técnica Base}
	
	\textbf{Receptor de imagen:} 35 $\times$ 43~cm. Orientación \textbf{longitudinal}
	(eje mayor del receptor paralelo al eje mayor de la región costal).
	
	\textbf{DFRI:} \textbf{100~cm} (1~metro). Distancia estándar para radiografía
	costal oblicua.
	
	\textbf{Técnica:} Básica o intermedia con \textbf{interés óseo}. Se utiliza
	\textbf{bajo kVp} para obtener alto contraste, permitiendo visualizar las
	costillas de extremo a extremo con su cortical ósea sin que el contraste del
	aire pulmonar enmascare posibles fisuras o fracturas.
	
	\textbf{Factores técnicos:} kVp \textbf{bajo} (técnica de alto contraste) ---
	mAs según espesor del hemitórax. El espesor se mide en la región de interés.
	
	\textbf{Bucky:} Bucky de pared con grilla. Se utiliza por el espesor de la
	región y para reducir la radiación dispersa.
	
	\subsubsection*{3. Posicionamiento}
	
	\paragraph{Paciente}
	De pie (bipedestación) en posición PA frente al Bucky de pared, o en decúbito
	semiprono si el paciente no puede estar de pie. El \textbf{lado afectado se
		aleja del receptor} a 45\textdegree. El paciente mira hacia adelante durante
	toda la exposición.
	
	\begin{figure}[h]
		\centering
		\includegraphics[width=0.7\linewidth]{imagenes/OA costillas}
		\caption{OAD}
		\label{fig:oa-costillas}
	\end{figure}
	
	\begin{figure}
		\centering
		\includegraphics[width=0.7\linewidth]{imagenes/OA1}
		\caption{}
		\label{fig:oa1}
	\end{figure}
	
	
	
	\paragraph{Región de interés}
	Rotación del cuerpo de \textbf{45\textdegree} respecto al plano del receptor:
	\begin{itemize}
		\item \textbf{OAI:} hombro anterior \textbf{izquierdo} apoyado contra el RI;
		el lado afectado (derecho) queda alejado a 45\textdegree. El brazo
		izquierdo (lado apoyado) con el codo flexionado y la mano sobre la
		cadera con la palma hacia afuera. El brazo derecho (lado alejado)
		elevado hasta la altura del hombro o tomado a la parte superior del
		Bucky para no tapar los pulmones.
		\item \textbf{OAD:} hombro anterior \textbf{derecho} apoyado contra el RI;
		el lado afectado (izquierdo) queda alejado a 45\textdegree. El brazo
		derecho (lado apoyado) con el codo flexionado y la mano sobre la
		cadera con la palma hacia afuera. El brazo izquierdo (lado alejado)
		elevado hasta la altura del hombro o tomado a la parte superior del
		Bucky.
	\end{itemize}
	
	\paragraph{Rayo Central}
	Perpendicular al plano del receptor. Punto de entrada: \textbf{T7}, equivalente
	a \textbf{8--10~cm por debajo de la escotadura yugular} (por debajo del nivel
	de la vértebra prominente C7).
	
	\paragraph{Receptor de imagen}
	El borde superior del receptor se coloca \textbf{4~cm por encima del borde
		superior del hombro}. La región a incluir abarca el \textbf{hemitórax de
		interés desde los vértices costales hasta los ángulos costofrénicos}.
	
	\paragraph{Respiración}
	\begin{itemize}
		\item \textbf{Costillas supra-diafragmáticas:} suspensión al final de una
		\textbf{inspiración profunda}, para descender el diafragma y liberar
		los arcos costales superiores.
		\item \textbf{Costillas infra-diafragmáticas:} suspensión al final de una
		\textbf{espiración profunda}, para elevar el diafragma y descubrir
		los arcos costales inferiores.
	\end{itemize}
	
	\subsubsection*{4. Anatomía Visible}
	
	Porción axilar de las costillas del hemitórax de interés desplegada y libre de
	superposición. Las costillas son visibles \textbf{a través del pulmón o del
		abdomen} según la región estudiada (supra o infra-diafragmática). La columna
	vertebral queda separada del borde lateral costal, generando una \textbf{doble
		distancia} entre la columna y el borde lateral de las costillas del lado
	afectado respecto al lado contralateral.
	
	\begin{itemize}
		\item \textbf{OAI:} despliega la curva axilar y los arcos costales anteriores
		\textbf{derechos}.
		\item \textbf{OAD:} visualiza las porciones \textbf{medias y anteriores del
			lado derecho} con mayor nitidez para la búsqueda de fracturas.
	\end{itemize}
	
	\subsubsection*{5. Criterios de Evaluación}
	
	\textbf{Posicionamiento correcto:} La porción axilar de las costillas del lado
	afectado se observa libre de superposición y desplegada. Distancia doble entre
	la columna vertebral y el borde lateral costal del lado afectado en comparación
	con el lado contralateral. Para costillas supra-diafragmáticas: se visualizan
	\textbf{8--9 arcos costales posteriores} por encima del diafragma y las
	costillas anteriores de la 5.ª a la 7.ª en su totalidad por encima del
	diafragma.
	
	\textbf{Ausencia de rotación:} A nivel de la columna, las bases de los cuerpos
	vertebrales presentan \textbf{un único borde}, lo que confirma que la rotación
	no supera los 45\textdegree. Si se visualiza doble borde en los cuerpos
	vertebrales, el paciente está rotado más de 45\textdegree (casi de perfil).
	
	\textbf{Penetración y contraste:} Las costillas se visualizan con su
	\textbf{cortical ósea nítida} a través del parénquima pulmonar o del abdomen.
	El alto contraste obtenido con bajo kVp permite identificar fisuras o fracturas
	sin que el contraste del aire enmascare los detalles óseos.
	
	\subsubsection*{6. Puntos Críticos}
	
	\textbf{Exceso de rotación ($>$45\textdegree):} El paciente queda casi de perfil.
	Radiológicamente se identifica porque las bases de los cuerpos vertebrales
	presentan \textbf{doble borde}. Corrección: verificar que el plano medio sagital
	forme exactamente 45\textdegree con el receptor antes de exponer.
	
	\textbf{Brazo del lado alejado mal posicionado:} El húmero y la escápula
	proyectan \textbf{demasiado hacia arriba} sobre el campo de interés porque el
	paciente no levantó el brazo ni lo alejó correctamente del cuerpo. Corrección:
	indicar al paciente que eleve el brazo hasta la altura del hombro o que lo
	apoye en la parte superior del Bucky de pared.
	
	\textbf{Adaptaciones:} En pacientes con trauma agudo que no toleran la
	bipedestación, realizar la proyección en \textbf{decúbito semiprono} con el
	mismo criterio de rotación de 45\textdegree. Ajustar la técnica si el paciente
	usa inmovilizador costal.
	
	\textbf{Trucos:}
	\begin{itemize}
		\item La nomenclatura de las oblicuas se da por el lado que \textbf{pega}
		contra el receptor, no por el lado afectado: recordar siempre que para
		ver un lado se apoya el \textbf{lado contrario}.
		\item Para la porción anterior y medial de los arcos costales, la oblicua
		estándar es insuficiente: se debe agregar una \textbf{segunda oblicua
			del mismo lado afectado} (acercando el lado afectado al receptor).
		\item Siempre anteceder las oblicuas con la \textbf{Rx de tórax PA a
			1{,}80~m} para descartar complicaciones pleuropulmonares antes de
		focalizar en la parrilla costal.
	\end{itemize}
	
\begin{figure*}
	\centering
	\includegraphics[width=0.7\linewidth]{imagenes/OA2}
	\caption{OAI - OPD}
	\label{fig:oa2}
\end{figure*}

\clearpage

%%%%%%%%%%%%%%%%%%%%%%%%%%%%%%%%%%%%%%%%%%%%%%%%%%%%%%%%%%%
%  SECCIÓN: Parrilla Costal — Oblicuas Posteriores
%  Archivo: secciones/torax.tex
%  Sin preámbulo ni \begin{document} — lo maneja main.tex
	%%%%%%%%%%%%%%%%%%%%%%%%%%%%%%%%%%%%%%%%%%%%%%%%%%%%%%%%%%%
	
	\section{Tórax}
	
	\subsection{Parrilla Costal - Oblicua Posterior AP - Porción Axilar (OPD / OPI)}
	
	\subsubsection*{1. Identificación}
	
	\textbf{Tipo:} Proyección complementaria. Forma par con las oblicuas anteriores (OAD / OAI) para completar el estudio de la parrilla costal.
	
	\textbf{Indicaciones clínicas principales:}
	\begin{itemize}
		\item Traumatismo costal con dolor axilar y/o posterior (ej.: caída de costado hacia atrás).
		\item Sospecha de fractura costal en arcos posteriores o borde axilar.
		\item Lesiones costales no visualizadas en la proyección PA estándar.
		\item Siempre complementar con radiografía de tórax PA estándar previa para descartar neumotórax, hemotórax o hemoneumotórax.
	\end{itemize}
	
	\subsubsection*{2. Técnica Base}
	
	\textbf{Receptor de imagen:} $35 \times 43$ cm. Orientación \textbf{longitudinal} (eje mayor vertical) en relación al eje mayor de la región.
	
	\textbf{DFRI:} 100 cm (1 metro).
	
	\textbf{Técnica:} Básica o intermedia para paciente estándar.
	
	\textbf{Factores técnicos:} kVp y mAs según técnica básica / intermedia según biotipo del paciente.
	
	\textbf{Bucky:} Sí, con grilla.
	
	\subsubsection*{3. Posicionamiento}
	
	\paragraph{Paciente}
	De pie (bipedestación) o en decúbito supino. El \textbf{lado afectado} se apoya contra el receptor de imagen. El plano medio sagital forma un ángulo de \textbf{45°} con respecto al plano del receptor.
	
	\paragraph{Región de interés}
	Se rota el cuerpo 45° sobre su eje longitudinal, de modo que el \textbf{hombro posterior del lado afectado} quede en contacto con el receptor:
	\begin{itemize}
		\item \textbf{OPD} (Oblicua Posterior Derecha): se acerca el lado derecho al receptor; el hombro posterior derecho apoya contra el RI. Se utiliza para visualizar los arcos costales posteriores \textbf{derechos}.
		\item \textbf{OPI} (Oblicua Posterior Izquierda): se acerca el lado izquierdo al receptor; el hombro posterior izquierdo apoya contra el RI. Se utiliza para visualizar los arcos costales posteriores \textbf{izquierdos} y el borde axilar izquierdo.
	\end{itemize}
	El \textbf{brazo del lado afectado} (más cercano al receptor) se levanta con la mano apoyada sobre la cabeza, para retirar la escápula del campo de interés. El \textbf{brazo contralateral} se abduce con la mano apoyada en la cadera.
	
	\begin{figure}[h]
		\centering
		\includegraphics[width=0.7\linewidth]{imagenes/OP1}
		\caption{}
		\label{fig:op1}
	\end{figure}
	
	\begin{figure}
		\centering
		\includegraphics[width=0.7\linewidth]{imagenes/OP3}
		\caption{OP}
		\label{fig:op3}
	\end{figure}
	
	
	\paragraph{Rayo Central}
	Perpendicular al plano del receptor. Punto de entrada: \textbf{T7} (a nivel de la séptima vértebra torácica), centrado en la mitad de la región de interés.
	
	\paragraph{Receptor de imagen}
	El borde superior del receptor se coloca \textbf{2 cm por encima de la vértebra prominente} (C7). De este modo se incluye la totalidad de los arcos costales de interés.
	
	\paragraph{Respiración}
	\begin{itemize}
		\item \textbf{Costillas supradiafragmáticas:} suspensión al final de una \textbf{inspiración profunda} (desciende el diafragma, se liberan las costillas superiores).
		\item \textbf{Costillas infradiafragmáticas:} suspensión al final de una \textbf{espiración profunda} (asciende el diafragma, permitiendo visualizar las costillas bajas sin superposición).
	\end{itemize}
	
	\subsubsection*{4. Anatomía Visible}
	
	Porción axilar de los arcos costales del lado afectado, libre de superposición y elongada. Arcos costales posteriores del hemitórax correspondiente (OPD para arcos derechos, OPI para arcos izquierdos). Las costillas deben ser visibles a través del parénquima pulmonar (región supradiafragmática) o a través del abdomen (región infradiafragmática) según la zona de estudio. Perfil de la tráquea rellena de aire como referencia de posicionamiento. Columna vertebral separada de la región de interés por la rotación de 45°.
	
	\subsubsection*{5. Criterios de Evaluación}
	
	\textbf{Posicionamiento correcto:} La porción axilar de las costillas del lado afectado aparece libre de superposición y alargada. La distancia entre la columna vertebral y el borde lateral de las costillas del lado afectado es aproximadamente el doble que la del lado contralateral. La escápula del lado afectado queda desplazada lateralmente, fuera del campo costal.
	
	\textbf{Ausencia de rotación excesiva:} Si las costillas aparecen acortadas con curvatura cerrada sobre sí mismas, el paciente está demasiado de perfil (más de 45°); se debe reducir la rotación. Si la columna superpone excesivamente el parénquima del lado afectado, la rotación también es excesiva.
	
	\textbf{Penetración y contraste:} Las costillas deben ser visibles tanto en la zona pulmonar como en la zona abdominal según la región estudiada. La técnica básica o intermedia debe garantizar visualización ósea con contraste adecuado.
	
	\subsubsection*{6. Puntos Críticos}
	
	\textbf{Rotación excesiva (más de 45°):} Las costillas se acortan y muestran curvatura cerrada, perdiendo el despliegue axilar buscado. Corrección: reducir el ángulo de rotación hasta 45° o menos.
	
	\textbf{Escápula superpuesta al borde costal externo:} Ocurre cuando el brazo del lado afectado no se levanta correctamente. Corrección: indicar al paciente que eleve el brazo con la mano sobre la cabeza para desplazar lateralmente la escápula fuera del campo.
	
	\textbf{Columna superpuesta al parénquima del lado afectado:} Indica rotación excesiva. Corrección: recolocar el paciente con rotación de 45° o menos.
	
	\textbf{Adaptaciones:} En pacientes con trauma severo o imposibilidad de bipedestación, la proyección puede realizarse en decúbito supino con el lado afectado elevado sobre el receptor. En ese caso, se mantiene la rotación de 45° con apoyo de cuñas o almohadillas radiolúcidas.
	
	\textbf{Trucos:}
	\begin{itemize}
		\item Regla mnemotécnica de nomenclatura: la oblicua se nombra según el lado que \textbf{pega} contra el receptor. OPD = lado derecho pegado al receptor.
		\item Para las oblicuas posteriores, siempre se \textbf{acerca el lado afectado}. Para las anteriores, se \textbf{aleja el lado afectado} (principio inverso).
		\item La OAD es equivalente a la OPI y viceversa (el cuerpo vertebral apunta hacia el mismo lado en ambas).
		\item Verificar la marca del lado en la imagen: debe coincidir con el lado de las apófisis espinosas visibles (OPD/OPI comparten este criterio con OAD/OAI respectivamente).
		\item Siempre comenzar el estudio costal con una PA de tórax a 1,80 m para descartar complicaciones pleurales antes de realizar las oblicuas.
	\end{itemize}
	
	\clearpage
	
	
	\begin{figure*}[h]
		\centering
		\includegraphics[width=0.7\linewidth]{imagenes/OP2}
		\caption{OPI - OAD}
		\label{fig:op2}
	\end{figure*}


\clearpage

	%----------------------------------------------------------
	%  Para agregar otro enfoque duplicá el bloque
	%  desde \subsection hasta \textbf{Trucos}
	%----------------------------------------------------------
	
\section{Tórax}

\subsection{Proyección Oblicua Anterior Derecha de Esternón (OAD)}

\subsubsection*{1. Identificación}

\textbf{Tipo:} Proyección de elección para el estudio del esternón. Complementa con lateral de esternón cuando se requiere.

\textbf{Indicaciones clínicas principales:}
\begin{itemize}
	\item Fractura esternal (contexto postraumático)
	\item Dolor retroesternal o dolor medio del pecho
	\item Procesos inflamatorios o infecciosos esternales
	\item Control de osteosíntesis esternal
\end{itemize}

\subsubsection*{2. Técnica Base}

\textbf{Receptor de imagen:} 24 × 30 cm. Orientación \textbf{longitudinal} en relación al eje mayor del esternón.

\textbf{DFRI:} 100 cm. Distancia estándar para proyecciones de tórax con Bucky mural.

\textbf{Técnica:} básica o intermedia. Se utiliza \textbf{kVp bajo} para obtener \textbf{alto contraste}, lo que permite resaltar la estructura ósea del esternón a través de la silueta cardiaca, aprovechando la densidad homogénea del mediastino como fondo.

\textbf{Factores técnicos:} kVp \textbf{bajo} --- mAs según espesor del tórax medido a nivel del esternón.

\textbf{Bucky:} sí, se utiliza Bucky con grilla.

\subsubsection*{3. Posicionamiento}

\paragraph{Paciente}
De pie (preferible) o en semidecúbito prono con ligera rotación. Brazo derecho al lado del cuerpo; brazo izquierdo elevado por encima de la cabeza.

\begin{figure}[h]
	\centering
	\includegraphics[width=0.7\linewidth]{imagenes/esternon1}
	\caption{}
	\label{fig:esternon1}
\end{figure}



\paragraph{Región de interés}
Paciente rotado \textbf{15 a 20°} hacia el lado derecho, alejando el lado izquierdo del receptor. La rotación varía según biotipo: tórax ancho requiere \textbf{menor} rotación; tórax delgado requiere \textbf{mayor} rotación. Para determinar el grado correcto, se puede colocar una mano sobre el esternón y otra sobre la apófisis espinosa correspondiente.

\paragraph{Rayo Central}
Perpendicular al plano del receptor. Punto de entrada: \textbf{centro del esternón}, a \textbf{2,5 cm a la izquierda de la línea media del paciente}.

\paragraph{Receptor de imagen}
Centrado en el esternón. El borde superior del receptor se ubica \textbf{4 cm por encima de la escotadura yugular} para asegurar la inclusión del manubrio.

\paragraph{Respiración}
Suspendida en \textbf{espiración}. La espiración eleva el diafragma y reduce el campo pulmonar superpuesto, mejorando el contraste sobre el esternón.

\begin{figure}[h]
	\centering
	\includegraphics[width=0.7\linewidth]{imagenes/esternon2}
	\caption{}
	\label{fig:esternon2}
\end{figure}


\subsubsection*{4. Anatomía Visible}

Estructuras principales: manubrio, cuerpo y apéndice xifoides del esternón; articulaciones esternoclaviculares; articulaciones costoesternales. El esternón debe proyectarse sobre la silueta cardiaca, libre de superposición de vértebras dorsales. La superposición sobre el mediastino es una ventaja: su densidad homogénea actúa como fondo uniforme que facilita la detección de alteraciones en el contorno y la cortical esternal.

\subsubsection*{5. Criterios de Evaluación}

\textbf{Posicionamiento correcto:} el esternón aparece superpuesto con la silueta cardiaca, sin contacto con las vértebras dorsales. Todo el eje longitudinal, desde el manubrio hasta el xifoides, debe estar visible en la imagen.

\textbf{Ausencia de rotación incorrecta:} el esternón no debe proyectarse sobre el campo pulmonar derecho; si se observa trama pulmonar superpuesta al hueso, la rotación es insuficiente o está invertida.

\textbf{Penetración y contraste:} técnica adecuada cuando el esternón es visible a través de las costillas, el corazón y el pulmón con alto contraste. Un kVp elevado produce escala larga que dificulta la diferenciación de bordes esternales y la detección de trazos de fractura.

\subsubsection*{6. Puntos Críticos}

\textbf{Proyección hacia el campo pulmonar derecho (OAI):} si el esternón se superpone con el parénquima pulmonar derecho, la trama vascular puede simular fracturas o lesiones óseas, con especial riesgo en pacientes añosos y fumadores. Siempre proyectar hacia la silueta cardiaca izquierda.

\textbf{Rotación insuficiente o excesiva:} la rotación insuficiente superpone el esternón sobre la columna vertebral; la excesiva lo desplaza hacia el pulmón derecho. El rango de \textbf{15 a 20°} debe ajustarse al biotipo.

\textbf{Adaptaciones:} para paciente politraumatizado que no puede oblicuarse, se realiza una \textbf{OPI en decúbito supino} angulando el RC \textbf{15 a 20° hacia el lado izquierdo}. La imagen tiene menor fidelidad diagnóstica; se valoran el contorno esternal y la integridad de la cortical.

\textbf{Trucos:} el esternón es un órgano hematopoyético y radiosensible; la OAD lo aleja de la fuente, reduciendo la dosis respecto a la OAI y produciendo menor magnificación geométrica.

\begin{figure*}[h]
	\centering
	\includegraphics[width=0.7\linewidth]{imagenes/esternon3}
	\caption{OAD esternón}
	\label{fig:esternon3}
\end{figure*}

\clearpage

%%%%%%%%%%%%%%%%%%%%%%%%%%%%%%%%%%%%%%%%%%%%%%%%%%%%%%%%%%%
%  SECCIÓN: Esternón
%  Archivo: secciones/esternon.tex
%  Sin preámbulo ni \begin{document} — lo maneja main.tex
	%%%%%%%%%%%%%%%%%%%%%%%%%%%%%%%%%%%%%%%%%%%%%%%%%%%%%%%%%%%
	
	\section{Tórax}
	
	\subsection{Proyección Lateral de Esternón}
	
	\subsubsection*{1. Identificación}
	
	\textbf{Tipo:} Proyección complementaria del esternón. Forma par con la OAD de esternón.
	
	\textbf{Indicaciones clínicas principales:}
	\begin{itemize}
		\item Fractura esternal (contexto postraumático)
		\item Dolor retroesternal o dolor medio del pecho
		\item Procesos inflamatorios o infecciosos esternales
		\item Control de osteosíntesis esternal
	\end{itemize}
	
	\subsubsection*{2. Técnica Base}
	
	\textbf{Receptor de imagen:} 24 × 30 cm. Orientación \textbf{longitudinal} en relación al eje mayor del esternón.
	
	\textbf{DFRI:} 140 a 180 cm. La distancia aumentada compensa la magnificación geométrica producida por la mayor distancia entre el esternón y el receptor en posición lateral.
	
	\textbf{Técnica:} intermedia. El esternón es un hueso plano, subcutáneo y rodeado de partes blandas y aire; se utiliza únicamente el espesor del esternón como referencia de medición. Se busca \textbf{máximo contraste óseo}.
	
	\textbf{Factores técnicos:} kVp \textbf{intermedio} --- mAs según espesor medido exclusivamente a nivel del esternón.
	
	\textbf{Bucky:} sí, se utiliza Bucky con grilla.
	
	\subsubsection*{3. Posicionamiento}
	
	\paragraph{Paciente}
	De pie (preferible) o en decúbito lateral. En bipedestación, el hombro del lado estudiado apoyado contra el receptor. Plano medio sagital paralelo al receptor de imagen.
	
	\paragraph{Región de interés}
	\textbf{En bipedestación --- posición de himno:} se instruye al paciente a sacar el pecho, llevar los brazos y hombros hacia atrás y entrelazar las manos por detrás de la espalda. Esta maniobra proyecta el esternón anteriormente, separándolo de las estructuras posteriores y reduciendo la superposición de costillas.
	
	\textbf{En decúbito lateral:} paciente recostado sobre un lado, brazos elevados por encima de la cabeza, hombros llevados hacia atrás. Posición lateral estricta, sin rotación.
	
	\paragraph{Rayo Central}
	Perpendicular al plano del receptor. Punto de entrada: \textbf{borde lateral del esternón a nivel de su zona media}.
	
	\paragraph{Receptor de imagen}
	Centrado en el esternón. El borde superior del receptor se ubica \textbf{4 cm por encima de la escotadura yugular} para incluir el manubrio. Colimar a los cuatro lados del área esternal.
	
	\paragraph{Respiración}
	Suspendida en \textbf{inspiración}. Al inspirar, los pulmones se llenan de aire y el esternón se desplaza anteriormente, mejorando su separación de las estructuras posteriores y aumentando el contraste.
	
	\subsubsection*{4. Anatomía Visible}
	
	Todo el esternón en perfil: manubrio, cuerpo y apéndice xifoides, desde la escotadura yugular hasta el extremo inferior del xifoides. El esternón debe visualizarse libre de superposición de costillas y con mínima superposición de partes blandas. No es una imagen de alta fidelidad diagnóstica; su valor principal reside en la evaluación del contorno anterior y posterior del hueso.
	
	\subsubsection*{5. Criterios de Evaluación}
	
	\textbf{Posicionamiento correcto:} todo el eje longitudinal del esternón visible desde el manubrio hasta el xifoides, sin superposición de costillas.
	
	\textbf{Ausencia de rotación:} el esternón se proyecta en perfil estricto; cualquier rotación del tronco superpone estructuras costales sobre el cuerpo esternal.
	
	\textbf{Penetración y contraste:} bordes óseos nítidos con máximo contraste. La nitidez de los bordes confirma además la ausencia de movimiento durante la exposición. Mínima superposición de partes blandas sobre la cortical esternal.
	
	\subsubsection*{6. Puntos Críticos}
	
	\textbf{Magnificación por distancia objeto-receptor:} en posición lateral el esternón queda alejado del receptor. Si no se compensa aumentando la DFRI a \textbf{140--180 cm}, se produce magnificación significativa que distorsiona la imagen y dificulta la evaluación de trazos de fractura.
	
	\textbf{Rotación del paciente:} cualquier rotación del tronco superpone las costillas sobre el esternón, oscureciendo el contorno óseo. La posición de himno en bipedestación es la principal herramienta para evitar este error.
	
	\textbf{Adaptaciones:} en paciente que no puede estar de pie se utiliza el \textbf{decúbito lateral} con brazos elevados y hombros hacia atrás. La calidad de imagen es comparable si se mantiene la posición lateral estricta.
	
	\textbf{Trucos:} medir el espesor únicamente a nivel del esternón y no del tórax completo para evitar sobreexposición. Esta proyección tiene limitaciones diagnósticas inherentes a su geometría y debe interpretarse siempre en conjunto con la OAD.
	
	%----------------------------------------------------------
	%  Para agregar otro enfoque duplicá el bloque
	%  desde \subsection hasta \textbf{Trucos}
	%----------------------------------------------------------